\section{Analytic Families}

RCOUNT packages provide clean, hash-bound inputs: source hashes, reporting-unit
lineage, normalized summaries, batch manifests, cast vote records, and audit
results. The W-series computes reproducible anomaly reports over those records.
The pipeline is uniform across families:

\begin{center}
certified/public results $\rightarrow$ normalized features $\rightarrow$
anomaly score $\rightarrow$ investigation queue.
\end{center}

The families differ in the question they ask and the RCOUNT records they
consume.

\begin{center}
\small
\begin{tabularx}{\textwidth}{lXl}
\toprule
Family & What it asks & RCOUNT input \\
\midrule
Turnout residuals & Is turnout unusual after accounting for history and demographics? & summaries, precinct lineage \\
Vote-share residuals & Are candidate shares unusual relative to comparable units? & summaries, CVRs \\
Batch/scanner effects & Do batches or scanners show systematic deviations? & batch manifests, CVRs \\
Digit tests & Do reported numbers show suspicious digit patterns? & summaries \\
Benford-style checks & Are leading-digit distributions unusual where applicable? & large count tables \\
Spatial outliers & Are neighboring units unexpectedly different? & reporting units, geography \\
Change-point tests & Did patterns shift abruptly across time or batches? & status events, batches \\
Audit discrepancy clustering & Do discrepancies cluster by source, machine, or batch? & audit observations \\
\bottomrule
\end{tabularx}
\end{center}

Each family is the subject of a later W-series paper. What unifies them is not a
single statistic but a shared input discipline and a shared output contract: the
features are derived from hash-bound RCOUNT records, and the result is an
anomaly report, never a pass or fail.
